% ECONOMICS_PROBLEM: {"id": "D82-308", "title": "Item pricing after optimal disclosure for two items", "jel_code": "D82", "source_id": "OP-0177"}
% !TeX program = xelatex
\documentclass[11pt,a4paper]{article}
\usepackage[margin=23mm]{geometry}
\usepackage{fontspec}
\setmainfont{Latin Modern Roman}
\usepackage{amsmath,amssymb,mathtools}
\usepackage{xcolor,enumitem,booktabs,longtable,array,xurl,hyperref}
\definecolor{muted}{HTML}{53636C}
\hypersetup{colorlinks=true,urlcolor=blue,linkcolor=blue}
\setlength{\parindent}{0pt}
\setlength{\parskip}{5pt}
\newcommand{\R}{\mathbb R}
\newcommand{\E}{\mathbb E}
\newcommand{\N}{\mathbb N}
\newcommand{\ind}{\mathbf 1}
\newcommand{\Var}{\operatorname{Var}}
\newcommand{\Cov}{\operatorname{Cov}}
\newcommand{\supp}{\operatorname{supp}}
\DeclareMathOperator*{\argmin}{arg\,min}
\DeclareMathOperator*{\argmax}{arg\,max}
\newcommand{\entrymeta}[3]{{\sffamily\small\color{muted}\textbf{\texttt{#1}}\quad|\quad #2\par #3\par}}
\newcommand{\Problemref}[1]{\texttt{#1}}
\newcommand{\JELref}[2]{\texttt{#2} (JEL \texttt{#1})}
\begin{document}
\section*{Item pricing after optimal disclosure for two items}
\textbf{Problem ID:} \texttt{D82-308}\quad \textbf{JEL:} \texttt{D82}\par
% BEGIN ECONOMICS STATEMENT
\label{problem:OP-0177}
% GAME_THEORY_PROBLEM: {"local_id": "GT-047", "original_number": 47, "kind": "P", "entry_kind": "statement_only", "published": false, "review_required": true, "category": "game_theory"}
\label{game:GT-047}
{\sffamily\small\color{muted}
\textbf{GT-047} | Information and mechanism design | \textbf{P}: explicit two-item conjecture.
\par}
\subsection*{Definitions and mathematical statement}
Let $F$ be a finitely supported distribution on $V\subset\mathbb R_{\geq0}^2$. The unit-demand buyer initially knows only $F$. The seller commits to a signal kernel $\sigma(ds\mid v)$, privately observed by the buyer, and a menu of lotteries. Write $\nu(s)=\mathbb E[v\mid s]$. A menu entry $(x,t)$ has $x\geq0$, $x_1+x_2\leq1$, $t\geq0$, and utility $\nu(s)\cdot x-t$. The menu includes $(0,0)$; consumption is chosen before values are realized. Let $R(F)$ be the supremum of expected payments over signal kernels and IC/IR menus. Let $R_{\rm item}(F)$ restrict menus to
\[
 \{(0,0),(e_1,p_1),(e_2,p_2)\},\qquad p_1,p_2\geq0,
\]
where $e_j$ assigns item $j$ with certainty. When utilities tie, select a maximizing entry with the highest payment. Determine whether
\[
 \forall F\text{ finitely supported on }\mathbb R_{\geq0}^2,\qquad
 R_{\rm item}(F)=R(F).
\]
IC means each signal type prefers its assigned entry to every other menu entry, and IR means that utility is nonnegative.
\subsection*{Short English statement}
Can suitable disclosure always make deterministic item pricing optimal when only two items are available?
\subsection*{Variants and scope boundaries}
Values may be correlated. The buyer does not know realized values before disclosure. Three or more items admit counterexamples; exogenous information and additive demand are different models.
\subsection*{Source relationship and status}
The two-item conjecture is explicit in the introduction and Section 3.3; Section 2 provides the model and seller-favoring tie convention. The finite-support statement avoids measure-theoretic extensions that are unnecessary to state the conjecture.
\subsection*{References}
\begingroup\footnotesize\raggedright
\textbf{[S1]} Yang Cai, Yingkai Li, and Jinzhao Wu (2025). \emph{Information Disclosure Makes Simple Mechanisms Competitive}. arXiv:2502.17809v1. \textit{Verified locator:} Section 1.1 after Example 3; Sections 2 and 3.3. \url{https://arxiv.org/html/2502.17809v1}
\par\endgroup

\subsection*{Common conventions: Source mathematical conventions}

\textbf{Finite games.} For a finite set $A$, $\Delta(A)$ is its probability
simplex. Players choose independent mixed strategies in Nash equilibrium;
correlation is allowed only when explicitly part of the solution concept.
In a game with payoff functions $u_i$ and strategy profile $x$, an additive
$\varepsilon$ Nash equilibrium satisfies
\[
 u_i(x)\geq u_i(x_i',x_{-i})-\varepsilon
 \quad\text{for every player $i$ and every admissible deviation $x_i'$}.
\]
For numerical approximation, payoffs are normalized as locally specified.
An $\varepsilon$ payoff residual is distinct from distance at most
$\varepsilon$ to the set of exact equilibria. Point convergence, convergence
of empirical play, and convergence in distance to an equilibrium set are
also distinct.

\textbf{Computational representations.} Unless stated otherwise, a complexity
question takes rational data with binary encoding; polynomial time is measured
in total bit length. A fixed-accuracy polynomial algorithm is not necessarily
polynomial in $1/\varepsilon$, and neither is necessarily polynomial in
$\log(1/\varepsilon)$. Succinct games and value, demand, or sampling oracles
must declare their interfaces. Exact real solutions require an explicit
representation; an unspecified list of arbitrary real numbers is not a
finite computational output. A claimed algorithmic barrier is meaningful
only for its specified model and reductions.

\textbf{Dynamic games.} Histories, information, observation, and allowable
randomization are part of the model. A stationary strategy depends only on
the current state; a positional strategy in a deterministic graph game
chooses one action at each controlled vertex. Nash equilibrium at an initial
state and equilibrium after every history are different requirements.
An expectation of a pathwise $\liminf$ payoff, a $\liminf$ of expected averages,
a limiting expected average, and a uniform finite-horizon equilibrium are
different criteria. None is silently substituted for another.

\textbf{Allocation and mechanisms.} A complete allocation assigns every item
exactly once unless otherwise stated. Zero-valued goods, negative values,
capacity constraints, feasibility, lotteries, and copies of goods affect
fairness definitions and are specified locally. Dominant-strategy incentive
compatibility, Bayesian incentive compatibility, universal truthfulness,
and truthfulness in expectation impose different inequalities. Whenever
an objective ratio has zero denominator, its convention is stated or those
instances are excluded.

\textbf{Infinite-dimensional models.} Expectations, integrals, stochastic
equations, and optimization objectives require their local regularity,
integrability, filtrations, and admissible domains. A backward--forward
mean-field system is a boundary-value problem; it is not an initial-value
semiflow without an additional construction. Long-time, large-population,
and vanishing-noise limits require an order of limits and a convergence
topology. If a minimum need not be attained, the objective is its infimum
or supremum and the approximation requirement is stated separately.
% END ECONOMICS STATEMENT
\end{document}
